% Subsection: Closure Under Arbitrary Unions  (notes stub; content authored later)
\subsection{Closure Under Arbitrary Unions}

\begin{remark*}[Exposition]
Arbitrary union closure is stronger than countable union closure because the
indexing set is unrestricted.  It is the union-side stability condition that
will later appear in topology.
\end{remark*}

\begin{definition}[Closed Under Arbitrary Unions]\label{def:closed-under-arbitrary-unions}
Let \(X\) be a set and let \(\mathcal{F}\subseteq\mathcal{P}(X)\).  The family
\(\mathcal{F}\) is \emph{closed under arbitrary unions} if for every indexed
family \(\{A_\alpha\}_{\alpha\in I}\) of members of \(\mathcal{F}\), where
\(I\) is any indexing set, one has
\[
  \bigcup_{\alpha\in I} A_\alpha \in \mathcal{F}.
\]
\end{definition}

\begin{remark*}[Standard quantified statement]
\[
  \forall I\,\forall \{A_\alpha\}_{\alpha\in I}\,
  \bigl((\forall \alpha\in I,\ A_\alpha\in\mathcal{F})
  \Rightarrow
  \bigcup_{\alpha\in I}A_\alpha\in\mathcal{F}\bigr).
\]
\end{remark*}

\begin{remark*}[Predicate reading]
\[
\begin{aligned}
&C_{\mathrm{arb}\cup}
=\mathsf{ClosureContext}(X,\mathcal{F}^{I},\mathcal{F},\bigcup_{I}),\\
&\operatorname{ClosedUnder}(\bigcup_{I},\mathcal{F},C_{\mathrm{arb}\cup}).
\end{aligned}
\]
\end{remark*}

\begin{remark*}[Interpretation]
Closure under arbitrary unions says that any collection of admitted subsets,
no matter how large its index set is, may be joined into one admitted subset.
\end{remark*}

\begin{remark*}[Examples]
The power set \(\mathcal{P}(X)\) is closed under arbitrary unions.  In later
chapters, a topology on \(X\) will be a set system whose open sets are closed
under arbitrary unions and finite intersections.
\end{remark*}

\begin{dependencies}
\begin{itemize}
  \item \hyperref[def:closed-under-operation]{Closed Under an Operation}
  \item \hyperref[def:arbitrary-operation-on-family]{Arbitrary Operation on a Family}
  \item \hyperref[def:indexed-family]{Indexed Family of Sets}
  \item \hyperref[def:indexed-union]{Indexed Union}
\end{itemize}
\end{dependencies}
